\section{Source reassociation with one-sided control bounds}
\label{sec:generalized-associativity}

The following ordinary mathematical proof uses exactly the structured
\texttt{run} semantics of Wade's register programs. It has not been
compiled in Lean or Rocq in this session.
Write $A;B$ for a cons cell with car $A$ and cdr $B$.

\paragraph{Reassociation theorem.}
Suppose no finite execution of $A$ exits with an escaping
$\operatorname{Break}(k)$, and no finite execution of $C$ exits with an
escaping $\operatorname{Continue}(k)$, for any positive $k$.
There is no restriction on $B$. Then, for every initial register vector
$R$, final vector $R'$, and exit $x$,
\[
 \operatorname{run}(A;(B;C),R,x,R')
 \quad\Longleftrightarrow\quad
 \operatorname{run}((A;B);C,R,x,R').
\]
In particular, $A$ may continue, $C$ may break, and any component may halt
or diverge. The assertion preserves the exact exit level and final
register vector, and assumes no total termination.

\paragraph{Finite-loop normal form.}
For a component $F$ and exit $x$, let $F_x$ be its binary relation on
register vectors. Relational composition $XY$ means first $X$, then $Y$;
$T^*$ is reflexive transitive closure by finite paths. Write the exits as
$N$, $H$, $\mathsf B_k$, and $\mathsf C_k$, with $k\geq 1$.
The constructors of \texttt{run} never produce a level-zero Break or
Continue exit, so these channels exhaust all possible finite exits.
For arbitrary $F,G$, put
\[
 T_{F,G}=F_NG_{\mathsf C_1}.
\]
This relation is exactly one restart of the cons cell. Define its
non-restarting alternatives by
\begin{align*}
 D_N^{F,G}&=F_{\mathsf B_1}\cup F_NG_N,\\
 D_{\mathsf B_k}^{F,G}
   &=F_{\mathsf B_{k+1}}\cup F_NG_{\mathsf B_k},\\
 D_{\mathsf C_k}^{F,G}
   &=F_{\mathsf C_k}\cup F_NG_{\mathsf C_{k+1}},\\
 D_H^{F,G}&=F_H\cup F_NG_H.
\end{align*}
The exact finite big-step relation is
\begin{equation}
 (F;G)_x=T_{F,G}^*D_x^{F,G}.
 \label{eq:cons-finite-loop-normal-form}
\end{equation}
To prove this identity, inspect the final constructor of a finite
\texttt{run} derivation. Every constructor except
\texttt{run\_node\_cdr\_continue1} is one of the displayed $D$ alternatives.
The exceptional constructor prepends one $T$ transition and has a
strictly smaller derivation for restarting the same cons cell. Induction
on the finite derivation extracts a finite $T$ path followed by $D$.
Conversely, a $D$ alternative gives a non-restart constructor, and
induction on the length of a finite $T$ path repeatedly applies the
restart constructor. This proves both inclusions.

\paragraph{Computing the left bracketing.}
The hypotheses say $A_{\mathsf B_k}=\varnothing$ and
$C_{\mathsf C_k}=\varnothing$ for every positive $k$. Put
$T=A_NB_{\mathsf C_1}$. Applying
\eqref{eq:cons-finite-loop-normal-form} gives
\begin{align*}
 (A;B)_N&=T^*A_NB_N,\\
 (A;B)_{\mathsf B_k}&=T^*A_NB_{\mathsf B_k},\\
 (A;B)_{\mathsf C_k}
   &=T^*(A_{\mathsf C_k}\cup A_NB_{\mathsf C_{k+1}}),\\
 (A;B)_H&=T^*(A_H\cup A_NB_H).
\end{align*}
The outer cell of $(A;B);C$ cannot restart because $C$ never continues.
Applying the normal form again and distributing composition over union,
we obtain
\begin{align}
 ((A;B);C)_N
   &=T^*(A_NB_{\mathsf B_1}\cup A_NB_NC_N),
   \label{eq:assoc-common-normal}\\
 ((A;B);C)_{\mathsf B_k}
   &=T^*(A_NB_{\mathsf B_{k+1}}\cup A_NB_NC_{\mathsf B_k}),
   \label{eq:assoc-common-break}\\
 ((A;B);C)_{\mathsf C_k}
   &=T^*(A_{\mathsf C_k}\cup A_NB_{\mathsf C_{k+1}}),
   \label{eq:assoc-common-continue}\\
 ((A;B);C)_H
   &=T^*(A_H\cup A_NB_H\cup A_NB_NC_H).
   \label{eq:assoc-common-halt}
\end{align}

\paragraph{Computing the right bracketing.}
The inner cell of $A;(B;C)$ cannot restart, so
\begin{align*}
 (B;C)_N&=B_{\mathsf B_1}\cup B_NC_N,\\
 (B;C)_{\mathsf B_k}&=B_{\mathsf B_{k+1}}\cup B_NC_{\mathsf B_k},\\
 (B;C)_{\mathsf C_k}&=B_{\mathsf C_k},\\
 (B;C)_H&=B_H\cup B_NC_H.
\end{align*}
Consequently the restart relation of $A;(B;C)$ is again exactly
$A_NB_{\mathsf C_1}=T$. Substitution in
\eqref{eq:cons-finite-loop-normal-form}, using
$A_{\mathsf B_k}=\varnothing$, gives precisely the four right-hand sides
of \eqref{eq:assoc-common-normal}--\eqref{eq:assoc-common-halt}.
All finite exit relations are equal, proving the theorem.

This proof includes arbitrarily many finite repetitions. A
$\operatorname{Continue}(1)$ emitted by $B$ repeats the same $A$-then-$B$
pair in both bracketings: the inner cell catches it in $(A;B);C$, and the
outer cell catches it in $A;(B;C)$. Infinite repetition or divergence
inside a component yields no additional finite derivation on either
side. No determinism hypothesis was used.

\paragraph{A computable sufficient guard.}
Let $(b_X,c_X)$ bound the escaping break and continue levels of a source
component $X$. The bounds used by the search satisfy
\[
 b_{X;Y}=\max\{\max(b_X-1,0),b_Y\},\qquad
 c_{X;Y}=\max\{c_X,\max(c_Y-1,0)\}.
\]
Increment and Halt have bounds $(0,0)$; Decrement has $(1,0)$;
Break$(k)$ has $(k,0)$; Continue$(k)$ has $(0,k)$.
Induction on finite \texttt{run} derivations shows that every emitted
positive level is at most its corresponding bound. The restart
constructor uses the induction hypothesis on its smaller restarted
derivation. Therefore a sufficient guard for this reassociation is just
\[
 b_A=0\quad\hbox{and}\quad c_C=0.
\]
The other four component bounds are unrestricted.

\paragraph{Application boundary.}
This is a source-program semantic theorem. A DAG rewrite must preserve
the component relations of $A,B,C$, or justify nested changes from the
inside out. A binary-table rewrite additionally needs a complete
simulation certificate accounting for every use of each changed row,
including backward-peek and runtime contexts. The theorem alone provides
neither a state-count saving nor correctness of a changed binary table.
